\section{The tensor-power formula and its virtual coefficients}

The formula on scan 28 is an identity of exact equivariant
Grothendieck groups.  An integral Cauchy filtration proves it on any
commutatively ringed topos, without a semisimplicity assumption.
Let $(X,\mathcal A)$ be such a topos, and let $K_0(X)$ be formed
from finite locally free $\mathcal A$-modules of bounded rank and
all short exact sequences between them.  Let $K_n(X)$ be the
analogous exact group with an $S_n$-action.  Induction makes
$\bigoplus_nK_n(X)$ a graded ring, on which $K_0(X)$ acts by
tensoring with modules with trivial action.

For a partition $\nu=(\nu_1,\ldots,\nu_l)$ of $n$, put
$\rho_\nu=\rho_{\nu_1}\cdots\rho_{\nu_l}$, and let $m_\nu$
be the monomial symmetric function, expressed as its integral
polynomial in elementary symmetric functions.  Evaluation at a
class of $K_0(X)$ uses its exterior powers.

\begin{theorem}\label{thm:tensor-formula}
For every $\xi\in K_0(X)$ one has
\[
 \tau^n(\xi)=\sum_{\nu\vdash n}m_\nu(\xi)\rho_\nu
 \quad\text{in }K_n(X).
\]
Here $\tau^n([E])=[E^{\otimes n}]$ with place permutations on
an effective class; the total series $\tau$ extends to virtual
classes by sending addition to the induction product.
\end{theorem}

\begin{proof}
We use the integral Schur functors $S_\lambda$ and the natural
Cauchy filtration
\[
 \operatorname{gr}\operatorname{Sym}^n(F\otimes G)
 \simeq\bigoplus_{\lambda\vdash n}S_\lambda(F)\otimes S_\lambda(G).
\]
Its existence over arbitrary commutative rings, and the freeness
and base-change compatibility of its factors, are the standard
basis and Cauchy theorems of
\citet[Theorem~II.2.16, p.~232; Theorem~III.1.4,
pp.~245--246]{ABW1982}.  Their notation has $L_{(a)}=\bigwedge^a$;
our convention is $S_\lambda=L_{\lambda'}$, where $\lambda'$
is the conjugate partition.  Its weight multiplicities are the
usual Kostka numbers $K_{\lambda\nu}$.
The construction consists of natural exterior diagonalizations
and symmetric multiplications.  The filtration and its free
quotients commute with scalar extension and changes of basis,
so they apply locally to free sheaves and glue on the topos.

Set $W=\Z^n$, with its diagonal torus and its permutation-matrix
$S_n$-action, and let
\[
 C_\lambda=(S_\lambda(W))_{(1,\ldots,1)}.
\]
This is a finite free integral $S_n$-lattice.  Weight extraction
is exact over $\Z$: a comodule for a split diagonal torus is a
direct sum of its weight modules, and all comodule maps preserve
that decomposition.  Thus extracting a weight from our filtration
preserves exactness and local freeness.

The $W$-weight $(1,\ldots,1)$ of
$\operatorname{Sym}^n(E\otimes W)$ is $E^{\otimes n}$, by the map
\[
 e_1\otimes\cdots\otimes e_n\longmapsto
 (e_1\otimes w_1)\cdots(e_n\otimes w_n).
\]
Permuting the $w_i$ corresponds to place permutations.  The first
weight extraction therefore gives the exact identity
\begin{equation}\label{eq:schur-tensor}
 [E^{\otimes n}]=\sum_{\lambda\vdash n}[S_\lambda(E)][C_\lambda].
\end{equation}
Apply this construction to $E=\Z^n$ with a second diagonal torus,
and take its weight $\nu$, padded by zeroes.  The tensor words of
that weight form the transitive Young coset set; the corresponding
weight space of $S_\lambda(E)$ has rank $K_{\lambda\nu}$.
Consequently, already over the integers,
\begin{equation}\label{eq:integral-Young}
 \rho_\nu=\sum_{\lambda\vdash n}K_{\lambda\nu}[C_\lambda].
\end{equation}
This equality follows from an exact filtration; equality of
rational characters is not being substituted for it.

For a third extraction, take the $W$-weight $\nu$ directly from
the Cauchy filtration.  Its left side is
$\bigotimes_j\operatorname{Sym}^{\nu_j}(E)$, whence
\[
 h_\nu(E)=\sum_{\lambda\vdash n}K_{\lambda\nu}[S_\lambda(E)].
\]
The identity $[\operatorname{Sym}^a(E)]=h_a(E)$ follows from
the exterior--symmetric Koszul complex, whose positive-degree
strands are exact locally on a free basis.  The Kostka matrix is
triangular with diagonal entries one in dominance order.  Comparing
the last display with the integral symmetric-function identity
$h_\nu=\sum_\lambda K_{\lambda\nu}s_\lambda$ and inverting this
matrix gives $[S_\lambda(E)]=s_\lambda(E)$ in $K_0(X)$.
Now $s_\lambda=\sum_\nu K_{\lambda\nu}m_\nu$, so
\eqref{eq:schur-tensor} and \eqref{eq:integral-Young} give
\[
 [E^{\otimes n}]
 =\sum_\nu m_\nu(E)\sum_\lambda K_{\lambda\nu}[C_\lambda]
 =\sum_\nu m_\nu(E)\rho_\nu.
\]

It remains to justify virtual classes.  For an exact sequence
$0\to E'\to E\to E''\to0$, the filtration of $E^{\otimes n}$
by the number of factors in $E'$ is $S_n$-stable and has factors
\[
 \operatorname{Ind}_{S_a\times S_{n-a}}^{S_n}
 \bigl((E')^{\otimes a}\otimes(E'')^{\otimes(n-a)}\bigr).
\]
Local splittings show that its submodules and quotients are
finite locally free.  Therefore $\tau(E)=\tau(E')\tau(E'')$
in the degree-completed induction ring.  Its constant coefficient
is one, so inversion of series extends $\tau$ to $K_0(X)$.
The symmetric-function identity
\[
 m_\nu(u+v)=\sum_{\alpha\cup\beta=\nu}m_\alpha(u)m_\beta(v)
\]
and $\rho_\alpha\rho_\beta=\rho_{\alpha\cup\beta}$ show that
the proposed right-hand total series has the same additivity
law.  These integral identities evaluate on $K_0(X)$, whose
special $\lambda$-ring structure is
\citeauthor{SGA6}'s theorem for commutatively ringed topoi
\citep[Expos\'e~VI, Theorem~3.3, pp.~392--393]{SGA6}.
The two series agree on effective classes and hence on all of
$K_0(X)$.
\end{proof}

For $n=2$, the formula is
\[
 [E^{\otimes2}]=\psi^2(E)\rho_2+
 [\bigwedge^2E]\rho_1^2,
 \qquad\psi^2(E)=[E]^2-2[\bigwedge^2E].
\]
The coefficients need not be effective.  For the standard
representation of $\mathrm{GL}_2$ in characteristic zero,
$\psi^2(E)=[\operatorname{Sym}^2E]-[\bigwedge^2E]$.
These are distinct irreducible representations, so the negative
coefficient precludes effectivity.  Thus the formula does not
give the direct-summand interpretation suggested in the margin.
